\documentclass[10pt,a4paper]{amsart}
\usepackage{amsmath,amssymb,amsthm,graphicx}
\usepackage{verbatim}
\newenvironment{reference}{\comment}{\endcomment}
%\newenvironment{reference}{}{}
\newenvironment{slogan}{\comment}{\endcomment}
\newenvironment{history}{\comment}{\endcomment}

% For commutative diagrams we use Xy-pic
\usepackage[all]{xy}

% We use 2cell for 2-commutative diagrams.
\xyoption{2cell}
% Theorem environments.
%
\theoremstyle{plain}
\newtheorem{theorem}[subsection]{Theorem}
\newtheorem{proposition}[subsection]{Proposition}
\newtheorem{lemma}[subsection]{Lemma}

\theoremstyle{definition}
\newtheorem{definition}[subsection]{Definition}
\newtheorem{example}[subsection]{Example}
\newtheorem{exercise}[subsection]{Exercise}
\newtheorem{situation}[subsection]{Situation}

\theoremstyle{remark}
\newtheorem{remark}[subsection]{Remark}
\newtheorem{remarks}[subsection]{Remarks}

\numberwithin{equation}{subsection}

% Macros
%
\def\lim{\mathop{\mathrm{lim}}\nolimits}
\def\colim{\mathop{\mathrm{colim}}\nolimits}
\def\Spec{\mathop{\mathrm{Spec}}}
\def\Hom{\mathop{\mathrm{Hom}}\nolimits}
\def\Ext{\mathop{\mathrm{Ext}}\nolimits}
\def\SheafHom{\mathop{\mathcal{H}\!\mathit{om}}\nolimits}
\def\SheafExt{\mathop{\mathcal{E}\!\mathit{xt}}\nolimits}
\def\Sch{\mathit{Sch}}
\def\Mor{\mathop{\mathrm{Mor}}\nolimits}
\def\Ob{\mathop{\mathrm{Ob}}\nolimits}
\def\Sh{\mathop{\mathit{Sh}}\nolimits}
\def\NL{\mathop{N\!L}\nolimits}
\def\CH{\mathop{\mathrm{CH}}\nolimits}
\def\proetale{{pro\text{-}\acute{e}tale}}
\def\etale{{\acute{e}tale}}
\def\QCoh{\mathit{QCoh}}
\def\Ker{\mathop{\mathrm{Ker}}}
\def\Im{\mathop{\mathrm{Im}}}
\def\Coker{\mathop{\mathrm{Coker}}}
\def\Coim{\mathop{\mathrm{Coim}}}

% Boxtimes
%
\DeclareMathSymbol{\boxtimes}{\mathbin}{AMSa}{"02}

%
% Macros for moduli stacks/spaces
%
\def\QCohstack{\mathcal{QC}\!\mathit{oh}}
\def\Cohstack{\mathcal{C}\!\mathit{oh}}
\def\Spacesstack{\mathcal{S}\!\mathit{paces}}
\def\Quotfunctor{\mathrm{Quot}}
\def\Hilbfunctor{\mathrm{Hilb}}
\def\Curvesstack{\mathcal{C}\!\mathit{urves}}
\def\Polarizedstack{\mathcal{P}\!\mathit{olarized}}
\def\Complexesstack{\mathcal{C}\!\mathit{omplexes}}
% \Pic is the operator that assigns to X its picard group, usage \Pic(X)
% \Picardstack_{X/B} denotes the Picard stack of X over B
% \Picardfunctor_{X/B} denotes the Picard functor of X over B
\def\Pic{\mathop{\mathrm{Pic}}\nolimits}
\def\Picardstack{\mathcal{P}\!\mathit{ic}}
\def\Picardfunctor{\mathrm{Pic}}
\def\Deformationcategory{\mathcal{D}\!\mathit{ef}}
\usepackage{cleveref}
\crefname{originallemma}{original lemma}{original lemmas}
\Crefname{originallemma}{Original lemma}{Original lemmas}
\makeatletter
\newlabel{lemma-site-abelian-category}{{05PM}{source}{}{sourcepointer-1}{}}
\newlabel{lemma-site-abelian-category@cref}{{[originallemma][0][]05PM}{[1][0][]source}{}{}{}}
\makeatother
\crefname{originalsection}{original section}{original sections}
\Crefname{originalsection}{Original section}{Original sections}
\makeatletter
\newlabel{homology-example-fibre-product-pushouts}{{05PJ}{source}{}{sourcepointer-2}{}}
\newlabel{homology-example-fibre-product-pushouts@cref}{{[originalsection][0][]05PJ}{[1][0][]source}{}{}{}}
\makeatother
\crefname{originallemma}{original lemma}{original lemmas}
\Crefname{originallemma}{Original lemma}{Original lemmas}
\makeatletter
\newlabel{homology-lemma-check-exactness}{{05AA}{source}{}{sourcepointer-3}{}}
\newlabel{homology-lemma-check-exactness@cref}{{[originallemma][0][]05AA}{[1][0][]source}{}{}{}}
\makeatother
\crefname{originallemma}{original lemma}{original lemmas}
\Crefname{originallemma}{Original lemma}{Original lemmas}
\makeatletter
\newlabel{categories-lemma-yoneda}{{001P}{source}{}{sourcepointer-4}{}}
\newlabel{categories-lemma-yoneda@cref}{{[originallemma][0][]001P}{[1][0][]source}{}{}{}}
\makeatother
\crefname{originallemma}{original lemma}{original lemmas}
\Crefname{originallemma}{Original lemma}{Original lemmas}
\makeatletter
\newlabel{homology-lemma-epimorphism-universal-abelian-category}{{05PK}{source}{}{sourcepointer-5}{}}
\newlabel{homology-lemma-epimorphism-universal-abelian-category@cref}{{[originallemma][0][]05PK}{[1][0][]source}{}{}{}}
\makeatother
\crefname{originallemma}{original lemma}{original lemmas}
\Crefname{originallemma}{Original lemma}{Original lemmas}
\makeatletter
\newlabel{sites-lemma-mono-epi-sheaves}{{00WN}{source}{}{sourcepointer-6}{}}
\newlabel{sites-lemma-mono-epi-sheaves@cref}{{[originallemma][0][]00WN}{[1][0][]source}{}{}{}}
\makeatother
\crefname{originalproposition}{original proposition}{original propositions}
\Crefname{originalproposition}{Original proposition}{Original propositions}
\makeatletter
\newlabel{sites-proposition-criterion-points}{{00YQ}{source}{}{sourcepointer-7}{}}
\newlabel{sites-proposition-criterion-points@cref}{{[originalproposition][0][]00YQ}{[1][0][]source}{}{}{}}
\makeatother
\title{Embedding an abelian category into the abelian sheaves on its associated site}
\author{The Stacks Project Authors}
\date{stacks-project:injectives.tex:local-snapshot-20261012; GNU FDL 1.2 or later}
\begin{document}
\maketitle
\noindent\textit{Editorial note.} All delivered mathematical text is the author's own, copied from the pinned Stacks Project chapter file sites-cohomology.tex; no statement or proof step was rewritten, reordered or removed.\par\medskip
\noindent\textit{Editorial note.} The proof is the author's complete argument for the exactness and full faithfulness, the only delegations being named results carried on this card as navigation stubs; the closing sentence about enough points names the criterion it applies and is preserved verbatim.\par\medskip
\noindent\textit{Editorial note.} Every labelled reference occurring inside the counted spans is carried in assembly.external\_references as a non-load-bearing navigation pointer (same-chapter ones at their definition site, cross-chapter ones at the invocation inside the counted span), so no load-bearing mathematics is replaced by a stub. No bibliography citation occurs inside the counted spans.\par\medskip
\noindent\textit{Editorial note.} The same-name check for this card scanned every staged Stacks chapter of this repository -- all sources/stacks-*/pinned-original-source.tex files, the card's own chapter excluded -- and the labels found are recorded in dedup.same\_name\_cross\_chapter on the card.\par\medskip
\noindent\textit{Editorial note.} The preamble is the chapter's own retained minimal preamble (sources/stacks-sites-cohomology/retained-preamble.tex), a verbatim excerpt of the native preamble: the Xy-pic setup and the native theorem declarations with the mathematical macros.\par\medskip
\noindent\textit{Allowed standard prerequisites (editorial).}\begin{itemize}
\item The site attached to an abelian category. Same chapter, Lemma lemma-site-abelian-category. Supplies the site structure used throughout, whose coverings make the representable presheaves sheaves.
\item Fibre products and the exactness criterion. Homology, Example homology-example-fibre-product-pushouts and Lemma homology-lemma-check-exactness. Supply the description of the fibre product of a surjection as a kernel and the criterion turning exactness into exactness of the induced sequences of Hom groups.
\item The universal property of epimorphisms. Homology, Lemma homology-lemma-epimorphism-universal-abelian-category. Supplies that the projection from the fibre product along a surjection is again a surjection, which produces the covering used for right exactness.
\item Monomorphisms and epimorphisms of sheaves. Sites, Lemma sites-lemma-mono-epi-sheaves and Proposition sites-proposition-criterion-points. Supply the characterization of surjectivity of a map of sheaves by local lifting and the criterion for a site to have enough points.
\item The Yoneda lemma. Categories, Lemma categories-lemma-yoneda. Supplies the full faithfulness of the Yoneda embedding used in the proof.
\end{itemize}\medskip
\ifdefined\hypertarget\hypertarget{sourcepointer-1}{}\fi
\noindent\textit{Source reference.} Original lemma 05PM: sources/stacks-injectives/pinned-original-source.tex, source lines 1267--1267. Navigation pointer invoked at proof line 1296; the counted span contains the pointer only and the cited result is not reproduced inside it, so the pointer is registered as a non-load-bearing navigation stub in the established form for Stacks material. No bibliography citation is involved.\par\medskip
\ifdefined\hypertarget\hypertarget{sourcepointer-2}{}\fi
\noindent\textit{Source reference.} Original section 05PJ: sources/stacks-homology/pinned-original-source.tex, source lines 918--918. Navigation pointer invoked at proof line 1308; the counted span contains the pointer only and the cited result is not reproduced inside it, so the pointer is registered as a non-load-bearing navigation stub in the established form for Stacks material. No bibliography citation is involved.\par\medskip
\ifdefined\hypertarget\hypertarget{sourcepointer-3}{}\fi
\noindent\textit{Source reference.} Original lemma 05AA: sources/stacks-homology/pinned-original-source.tex, source lines 967--967. Navigation pointer invoked at proof line 1328; the counted span contains the pointer only and the cited result is not reproduced inside it, so the pointer is registered as a non-load-bearing navigation stub in the established form for Stacks material. No bibliography citation is involved.\par\medskip
\ifdefined\hypertarget\hypertarget{sourcepointer-4}{}\fi
\noindent\textit{Source reference.} Original lemma 001P: sources/stacks-categories/pinned-original-source.tex, source lines 758--758. Navigation pointer invoked at proof line 1334; the counted span contains the pointer only and the cited result is not reproduced inside it, so the pointer is registered as a non-load-bearing navigation stub in the established form for Stacks material. No bibliography citation is involved.\par\medskip
\ifdefined\hypertarget\hypertarget{sourcepointer-5}{}\fi
\noindent\textit{Source reference.} Original lemma 05PK: sources/stacks-homology/pinned-original-source.tex, source lines 1136--1136. Navigation pointer invoked at proof line 1283; the counted span contains the pointer only and the cited result is not reproduced inside it, so the pointer is registered as a non-load-bearing navigation stub in the established form for Stacks material. No bibliography citation is involved.\par\medskip
\ifdefined\hypertarget\hypertarget{sourcepointer-6}{}\fi
\noindent\textit{Source reference.} Original lemma 00WN: sources/stacks-sites/pinned-original-source.tex, source lines 2473--2473. Navigation pointer invoked at proof line 1346; the counted span contains the pointer only and the cited result is not reproduced inside it, so the pointer is registered as a non-load-bearing navigation stub in the established form for Stacks material. No bibliography citation is involved.\par\medskip
\ifdefined\hypertarget\hypertarget{sourcepointer-7}{}\fi
\noindent\textit{Source reference.} Original proposition 00YQ: sources/stacks-sites/pinned-original-source.tex, source lines 9618--9618. Navigation pointer invoked at proof line 1350; the counted span contains the pointer only and the cited result is not reproduced inside it, so the pointer is registered as a non-load-bearing navigation stub in the established form for Stacks material. No bibliography citation is involved.\par\medskip
\begin{lemma}
\label{lemma-embedding}
Let $\mathcal{A}$ be an abelian category.
Let $\mathcal{C} = (\mathcal{A}, \text{Cov})$ be the
site defined in
Lemma \ref{lemma-site-abelian-category}.
Then $X \mapsto h_X$ defines a fully faithful, exact functor
$$
\mathcal{A} \longrightarrow \textit{Ab}(\mathcal{C}).
$$
Moreover, the site $\mathcal{C}$ has enough points.
\end{lemma}

\begin{proof}
Suppose that $f : V \to U$ is a surjective morphism of $\mathcal{A}$.
Let $K = \Ker(f)$. Recall that
$V \times_U V = \Ker((f, -f) : V \oplus V \to U)$, see
Homology, Example \ref{homology-example-fibre-product-pushouts}.
In particular there exists an injection $K \oplus K \to V \times_U V$.
Let $p, q : V \times_U V \to V$ be the two projection morphisms.
Note that $p - q : V \times_U V \to V$ is a morphism such that
$f \circ (p  - q) = 0$. Hence $p - q$ factors through $K \to V$.
Let us denote this morphism by $c : V \times_U V \to K$.
And since the composition $K \oplus K \to V \times_U V \to K$
is surjective, we conclude that $c$ is surjective. It follows that
$$
V \times_U V \xrightarrow{p - q} V \to U \to 0
$$
is an exact sequence of $\mathcal{A}$.
Hence for an object $X$ of $\mathcal{A}$ the sequence
$$
0 \to
\Hom_\mathcal{A}(U, X) \to
\Hom_\mathcal{A}(V, X) \to
\Hom_\mathcal{A}(V \times_U V, X)
$$
is an exact sequence of abelian groups, see
Homology, Lemma \ref{homology-lemma-check-exactness}.
This means that $h_X$ satisfies the sheaf condition
on $\mathcal{C}$.

\medskip\noindent
The functor is fully faithful by
Categories, Lemma \ref{categories-lemma-yoneda}.
The functor is a left exact functor between abelian categories by
Homology, Lemma \ref{homology-lemma-check-exactness}.
To show that it is right exact, let $X \to Y$ be a surjective morphism
of $\mathcal{A}$. Let $U$ be an object of $\mathcal{A}$, and let
$s \in h_Y(U) = \Mor_\mathcal{A}(U, Y)$ be a section of $h_Y$
over $U$. By
Homology, Lemma \ref{homology-lemma-epimorphism-universal-abelian-category}
the projection $U \times_Y X \to U$ is surjective.
Hence $\{V = U \times_Y X \to U\}$ is a covering of $U$ such that
$s|_V$ lifts to a section of $h_X$. This proves that
$h_X \to h_Y$ is a surjection of abelian sheaves, see
Sites, Lemma \ref{sites-lemma-mono-epi-sheaves}.

\medskip\noindent
The site $\mathcal{C}$ has enough points by
Sites, Proposition \ref{sites-proposition-criterion-points}.
\end{proof}

\end{document}
